% Lettre de motivation (FR) — El Mehdi Haress
% Cible : Analyste Quantitatif Front Office (H/F) — Canopee, Paris
% Compilation : xelatex lettre-canopee-quant.tex
\documentclass[10pt,a4paper]{article}

\usepackage{fontspec}
\usepackage[french]{babel}
\usepackage[T1]{fontenc}
\usepackage[
  top=1.5cm, bottom=1.3cm, left=2.0cm, right=2.0cm
]{geometry}
\usepackage{xcolor}
\usepackage{microtype}
\usepackage[hidelinks]{hyperref}
\usepackage{fontawesome5}

% ---------- Polices (identiques au CV) ----------
\setmainfont{texgyrepagella}[
  Extension      = .otf,
  Scale          = 1.05,
  UprightFont    = *-regular,
  BoldFont       = *-bold,
  ItalicFont     = *-italic,
  BoldItalicFont = *-bolditalic,
]
\newfontfamily\headingfont{texgyreheros}[
  Extension      = .otf,
  UprightFont    = *-regular,
  BoldFont       = *-bold,
  ItalicFont     = *-italic,
  BoldItalicFont = *-bolditalic,
]

% ---------- Couleurs (identiques au CV) ----------
\definecolor{accent}{RGB}{28,58,94}
\definecolor{ink}{RGB}{26,26,26}
\definecolor{muted}{RGB}{95,105,120}
\definecolor{rulegray}{RGB}{200,206,214}
\color{ink}

% Caractères laissés tels quels : un kern négatif casse l'extraction du mot-clé « C++ » par les ATS.
\newcommand{\cpp}{\mbox{C++}}

% ---------- Mise en page ----------
\pagestyle{empty}
\setlength{\parindent}{0pt}
\setlength{\parskip}{1.25ex plus 0.6ex minus 0.15ex}
\linespread{1.10}
\frenchspacing

\begin{document}

% =====================  EN-TÊTE  =====================
{\headingfont\bfseries\fontsize{20}{23}\selectfont EL MEHDI HARESS}\\[0.7ex]
{\headingfont\fontsize{10.5}{13}\selectfont\color{accent}%
Docteur en mathématiques appliquées \textbar{} Ingénieur CentraleSupélec}\\[1.0ex]
{\small
\faEnvelope[regular]~\href{mailto:el-mehdi.haress@inria.fr}{el-mehdi.haress@inria.fr}\quad
\faPhone*~+33 7 58 23 41 50\quad
\faLinkedin~\href{https://www.linkedin.com/in/el-mehdi-haress-67750b3aa/}{el-mehdi-haress}\quad
\faGlobe~\href{https://elmehdiharess.org}{elmehdiharess.org}
}

\vspace{0.4ex}
{\color{rulegray}\rule{\linewidth}{0.7pt}}

\vspace{1.2ex}

% =====================  DESTINATAIRE  =====================
\begin{minipage}[t]{0.5\textwidth}
\vspace{0pt}\raggedright\small
\textbf{Canopee}\\
Practice Quant \textemdash{} Finance quantitative\\
77 boulevard Berthier, 75017 Paris
\end{minipage}\hfill
\begin{minipage}[t]{0.42\textwidth}
\vspace{0pt}\raggedleft\small
Sophia Antipolis, le 16 septembre 2026
\end{minipage}

\vspace{2.2ex}

{\small\textbf{Objet :} candidature au poste d'Analyste Quantitatif Front Office (H/F) \textemdash{} Paris}

\vspace{2.2ex}

Madame, Monsieur,

La plupart des offres de conseil vendent une mission. La vôtre vend aussi une \textbf{Practice Quant} : des
travaux de recherche \og~innovants et alternatifs~\fg{} en \textbf{modélisation stochastique} et en
\textbf{valorisation d'instruments exotiques}. C'est cette phrase, plus que le titre, qui m'a fait
candidate. Je viens de cinq ans de recherche sur ces objets ; je cherche maintenant le book, le produit et
l'échéance qui les rendent concrets.

Mes travaux portent sur les \textbf{équations différentielles stochastiques dirigées par un mouvement brownien
fractionnaire} : dynamiques à mémoire longue, trajectoires irrégulières, calibration délicate. Ce sont les
objets sous-jacents de la \textbf{volatilité rough} \textemdash{} rough Heston, rough Bergomi \textemdash{} et,
plus généralement, de la famille de modèles que vous citez : \textbf{volatilité stochastique, locale, à
sauts}. L'un de mes articles construit des estimateurs conjoints de la dérive, de la volatilité et de
\textbf{l'indice de Hurst} sur observations discrètes, avec vitesses de convergence prouvées. C'est, formulé
en chercheur, exactement la question d'une calibration de front office :
\emph{cet estimateur converge-t-il, à quelle vitesse, et sous quelles hypothèses se dégrade-t-il ?}

Ma thèse prolonge ce réflexe côté \textbf{schémas numériques}. Elle établit un cadre d'approximation pour des
dynamiques \textbf{à dérive singulière}, là où les hypothèses habituelles des schémas d'Euler tombent. À
Leeds, j'ai ensuite quantifié la \textbf{sensibilité d'une solution aux perturbations de sa spécification et
de ses paramètres} : dans un langage de trading, ce sont les \textbf{Grecs} et la robustesse du pricer. J'ai
enseigné la \textbf{finance stochastique et la modélisation du risque en M2 à CentraleSupélec} : expliquer un
modèle, en assumer les limites, et former ceux qui devront s'en servir \textemdash{} y compris, demain, les
utilisateurs d'une librairie de pricing.

Ce qui m'attire chez Canopee est le double terrain. Côté client, développer des méthodes, les calibrer, les
intégrer dans une librairie et les mettre entre les mains des traders. Côté Practice, continuer à chercher,
sur des modèles moins sages que Heston, sans attendre que la littérature soit devenue un standard de marché.
Je code en \textbf{Python} au quotidien depuis huit ans, j'ai les bases du \cpp{} de ma formation d'ingénieur,
et j'ai livré seul plusieurs applications jusqu'à la production, pour un client réel : monter en compétence
sur le C\# et contribuer à une librairie de pricing ne me fait pas peur. Je suis mobile sur l'Île-de-France.

Je serais très heureux de vous montrer, de vive voix, ce que cinq ans de stochastique rough et singulière
peuvent apporter à un book de front office.

Veuillez agréer, Madame, Monsieur, l'expression de mes salutations distinguées.

\vspace{2.6ex}

\hfill
\begin{minipage}{0.42\textwidth}
\raggedleft
{\headingfont\bfseries\small El Mehdi Haress}\\[0.2ex]
{\small\color{muted}Docteur en mathématiques appliquées}
\end{minipage}

\end{document}
